% ==========================
% Chapter 1
% ==========================
\CustomChapter{廣義相對論之理論框架}
\label{ch:chap1}
\ChapterQuote{想象远比知识重要}{阿尔伯特 爱因斯坦}

\section{緒論}
% This chapter introduces the theoretical foundations required for the study developed in this thesis. We first recall the geometric formulation of general relativity through the Einstein field equations, then present the Schwarzschild solution as the reference exact solution used throughout this work \cite{einstein1916,wald1984}.
本章介紹本論文所進行之研究所需的理論基礎。我們首先透過愛因斯坦場方程回顧一般相對論的幾何表述，接著提出施瓦西解，作為本研究中全程使用的參考精確解\cite{einstein1916,wald1984}。

% ==========================
\section{幾何基礎}
% ==========================

% General relativity describes gravitation not as a force, but as a manifestation of spacetime curvature. The geometry of spacetime is encoded in the metric tensor $g_{\mu\nu}$, which defines the invariant interval
廣義相對論將引力描述為時空彎曲的表現形式，而非一種力。時空的幾何結構蘊含於度量張量 $g_{\mu\nu}$ 中，該張量定義了不變區間
\begin{equation}
	ds^2 = g_{\mu\nu}\,dx^{\mu}dx^{\nu},
	\label{eq:metric}
\end{equation}
using the Einstein summation convention over repeated indices \cite{misner1973}.\marginpar{\footnotesize Greek indices run from $0$ to $3$, with $x^0=ct$.}

The curvature of spacetime is related to the distribution of matter and energy through the Einstein field equations
\begin{equation}
	G_{\mu\nu} \;=\; R_{\mu\nu} - \tfrac{1}{2}R\,g_{\mu\nu} \;=\; \frac{8\pi G}{c^4}\,T_{\mu\nu},
	\label{eq:efe}
\end{equation}
where $R_{\mu\nu}$ is the Ricci tensor, $R$ the scalar curvature, and $T_{\mu\nu}$ the stress-energy tensor \cite{wald1984}. As stated in Eq.~\eqref{eq:efe}, the left-hand side is purely geometric while the right-hand side encodes the physical matter content, illustrating the interplay expressed by Wheeler's summary: matter tells spacetime how to curve, and spacetime tells matter how to move\footnote{This heuristic formulation, attributed to J. A. Wheeler, is widely used in introductory treatments of general relativity \cite{misner1973}.}.

% ==========================
\section{The Schwarzschild Solution}
% ==========================

For a static, spherically symmetric mass distribution in vacuum ($T_{\mu\nu}=0$), Eq.~\eqref{eq:efe} admits the exact Schwarzschild solution \cite{schwarzschild1916}
\begin{equation}
	ds^2 = -\left(1-\frac{2GM}{c^2 r}\right)c^2 dt^2 + \left(1-\frac{2GM}{c^2 r}\right)^{-1} dr^2 + r^2\, d\Omega^2,
	\label{eq:schwarzschild}
\end{equation}
where $d\Omega^2 = d\theta^2 + \sin^2\theta\, d\phi^2$ is the metric on the unit sphere. The quantity
\begin{equation}
	r_s = \frac{2GM}{c^2}
	\label{eq:rs}
\end{equation}
defines the Schwarzschild radius, corresponding to the event horizon of a non-rotating black hole \cite{chandrasekhar1983}. Equation~\eqref{eq:schwarzschild} reduces to the flat Minkowski metric as $r \to \infty$, recovering special relativity far from the mass \cite{hartle2003}.

\begin{figure}[H]
	\centering
	\includegraphics[width=0.55\textwidth]{gfx/figure1}
	\caption{Embedding diagram illustrating the curvature of spacetime around a spherical mass, as described by the Schwarzschild metric (Eq.~\eqref{eq:schwarzschild}).}
	\label{fig:schwarzschild}
\end{figure}

% ==========================
\section{Comparison of Exact Solutions}
% ==========================

Table~\ref{tab:solutions} summarizes the main properties of the exact solutions considered as reference cases in relativistic astrophysics.

\begin{table}[H]
	\centering
	\caption{Comparison of classical exact solutions of the Einstein field equations.}
	\label{tab:solutions}
	\begin{tabular}{lccc}
		\hline
		\textbf{Solution} & \textbf{Symmetry} & \textbf{Charge} & \textbf{Angular momentum} \\
		\hline
		Schwarzschild \cite{schwarzschild1916} & Spherical & No & No \\
		Reissner--Nordstr\"om \cite{reissner1916} & Spherical & Yes & No \\
		Kerr \cite{kerr1963} & Axial & No & Yes \\
		Kerr--Newman \cite{newman1965} & Axial & Yes & Yes \\
		\hline
	\end{tabular}
\end{table}

As shown in Table~\ref{tab:solutions}, the Schwarzschild metric of Eq.~\eqref{eq:schwarzschild} corresponds to the simplest case, from which the more general rotating and charged solutions can be derived \cite{kerr1963}.

% ==========================
\section{Conclusion}
% ==========================

This chapter recalled the geometric foundations of general relativity, culminating in the derivation of the Schwarzschild solution (Eq.~\eqref{eq:schwarzschild}), which serves as the reference framework for the analysis developed in the remainder of this thesis \cite{einstein1916,wald1984}.